\documentclass[11pt]{article}

\usepackage[margin=1in]{geometry}
\usepackage{amsmath,amssymb,booktabs,longtable,array}
\usepackage{enumitem}
\usepackage[hidelinks]{hyperref}

\title{Real-Ecology Benchmark Algorithms:\\Method Equations, Training Flow, and Implementation Adaptations}
\author{For readers who know the algorithm papers but do not inspect the code}
\date{July 2026}

\newcommand{\code}[1]{\texttt{#1}}
\newcommand{\E}{\mathbb{E}}
\newcommand{\Var}{\mathrm{Var}}
\newcommand{\1}{\mathbf{1}}

\begin{document}
\sloppy
\maketitle

\section{Purpose}

This document explains how the benchmark algorithms are instantiated in the
real-ecology continuous-observation setting.  It is written for a reader who
knows the relevant algorithm papers or method families, but does not inspect the
source code.  The goal is to describe:
\begin{itemize}[leftmargin=*]
  \item the mathematical form of the benchmark POMDP;
  \item the detailed training and decision flow for each algorithm;
  \item the update equations used by each adaptation;
  \item how each method is evaluated;
  \item what was changed from the original paper-style method and why that
  adaptation was needed for this benchmark.
\end{itemize}

The remainder of the note presents the implementation design and the intended
interpretation of each method.

\section{Benchmark POMDP Shared by All Methods}

\subsection{State, Observation, and Public Belief}

The latent ecological state is continuous abundance $s_t\ge 0$.  The agent does
not observe $s_t$ directly.  It observes a noisy survey:
\[
  o_t = s_t\epsilon_t,\qquad
  \epsilon_t\sim\mathrm{LogNormal}(0,\sigma_{\mathrm{obs}}).
\]
The decision input to ordinary methods is therefore a public belief over
abundance and latent structural variables, constructed from the public history:
\[
  h_t=(o_0,a_0,o_1,\ldots,a_{t-1},o_t)
\]
plus public control accumulators.  The internal belief representation is a
particle approximation:
\[
  b_t \approx \{(s_t^{(i)}, z_t^{(i)}, g_t^{(i)}, w_t^{(i)})\}_{i=1}^N,
\]
where $z_t$ denotes continuous latent context particles, $g_t$ denotes regime
particles, and $w_t$ are normalized particle weights.

Realized rewards are stored in the offline dataset for supervised value targets
and evaluation summaries, but they are not part of the online
\code{PublicTransition} passed to \code{policy.observe()}.  Online filtering and
posterior updates therefore consume only the next observation, done flags, and
public control fields.  This is the current leakage boundary in the
implementation.

\subsection{Real-Ecology Actions}

The benchmark uses 11 real-data actions loaded from the single authoritative
table directory \code{discrete\_action\_cont\_obser/real\_ecology\_data}.  The
loader validates 9 populations, 11 actions, and every population/action effect
row.  Each action $a$ for population $p$ and family $f$ has:
\[
  \Delta r(a,p,f),\qquad \Delta K(a,p),\qquad
  \Delta N(a,p),\qquad c(a,p).
\]
The public cumulative controls evolve as:
\begin{align}
  \rho_{t+1} &= (1-d_r)\rho_t+\Delta r(a_t,p,f),\\
  \kappa_{t+1} &= (1-d_K)\kappa_t+\Delta K(a_t,p),\\
  K^{\mathrm{eff}}_{t+1}
    &= \mathrm{clip}\{K_{\mathrm{base}}+\kappa_{t+1},K_{\min},K_{\max}\}.
\end{align}
For the real setting, $d_r=1$ and $d_K=0$, so $\rho_{t+1}$ is the current action's
set-point growth/mortality value while $\kappa_t$ accumulates habitat capacity
improvements.

The real ecological growth rate is:
\[
  r^{\mathrm{eff}}_{t+1}
  = \mathrm{clip}\{\rho_{t+1},r_{\min}(p),r_{\max}(p)\}.
\]
The hidden baseline growth draw is still present in the shared environment
interface, but in \code{control\_mode="real\_setpoint"} it is not added to the
set-point; the action table supplies the effective set-point axis.
Direct translocation enters as:
\[
  m_t = \max\{s_t+\Delta N(a_t,p),0\}.
\]
Only the translocation action has nonzero $\Delta N$.

\subsection{Ecological Transition}

Negative effective growth is treated as unconditional mortality rather than as a
negative coefficient inside the density term:
\[
  r_+ = \max(r^{\mathrm{eff}},0),\qquad
  r_- = \min(r^{\mathrm{eff}},0).
\]
For the Ricker family:
\[
  s_{t+1}
  =
  m_t\exp\left(r_+\left(1-\frac{m_t}{K^{\mathrm{eff}}_{t+1}}\right)\right)
  \exp(r_-).
\]
For the Allee family:
\[
  s_{t+1}
  =
  m_t\exp\left(
    r_+\left(1-\frac{m_t}{K^{\mathrm{eff}}_{t+1}}\right)
       \left(\frac{m_t}{C}-1\right)
  \right)\exp(r_-).
\]
For the theta-logistic family:
\[
  s_{t+1}
  =
  \left[m_t+r_+m_t
  \left(1-\left(\frac{m_t}{K^{\mathrm{eff}}_{t+1}}\right)^\theta\right)\right]
  \exp(r_-).
\]
The regime family switches between weak and strong Allee-like maps with a
private regime variable.

\subsection{Reward}

The reward is state-dependent and uses the true next abundance:
\[
  R_t =
  \alpha\frac{s_{t+1}}{s_{t+1}+K_{\mathrm{ref}}(p)}
  - c(a_t)
  - P_m\chi_t.
\]
For real cells $K_{\mathrm{ref}}(p)=K_{\mathrm{base}}(p)$, so the benefit is
anchored to the published baseline capacity rather than to the improved
$K^{\mathrm{eff}}_t$.
The two reward modes share the same benefit and cost.  They differ only in the
collapse penalty:
\[
  P_{\mathrm{yield}}=0,\qquad
  P_{\mathrm{safe}}=\text{configured collapse penalty}.
\]
The default real safe-mode indicator is occupancy:
\[
  \chi_t=\1\{s_{t+1}\le s_{\mathrm{safe}}(p)\}.
\]
This is deliberate: depleted populations can start below the safety floor, so a
downward-crossing-only penalty would not charge them while they remain unsafe.
The old event-only convention is still available as an explicit ablation:
\[
  \chi_t=\1\{s_t>s_{\mathrm{safe}}(p),\ s_{t+1}\le s_{\mathrm{safe}}(p)\}.
\]
Safety floors are depletion-aware by default: healthy large populations use
$0.10K_{\mathrm{base}}$, mid-scale or moderately depleted populations use
$0.20K_{\mathrm{base}}$, and the smallest or most depleted cases use
$0.25K_{\mathrm{base}}$.  The separate absolute MVP diagnostic threshold is
50 individuals and is recorded for evaluation, not used as the reward floor.

\subsection{Shared Evaluation Battery}

Every trained policy is evaluated by rolling out from $s_0=N_0(p)$ under the
same observation model and seed protocol.  The evaluation records operational
return, true return, collapse entry, unsafe fraction, MVP fraction and breach,
final abundance, minimum abundance, persistence, economic cost, danger-zone
action fractions, filter errors, runtime, backend metadata, and fallback counts.
Reward-mode-specific policies are trained separately.

\section{Shared Training Infrastructure}

\subsection{Offline Dataset and Belief Cache}

The offline public dataset contains training transitions:
\[
  (o_t,a_t,R_t,o_{t+1},d_t,\rho_t,\kappa_t,K^{\mathrm{eff}}_t,
   \rho_{t+1},\kappa_{t+1},K^{\mathrm{eff}}_{t+1}).
\]
The reward $R_t$ in this tuple is a logged training target, not an online
observation supplied to deployed policies.  Private states and hidden parameters
are saved in a separate sidecar and used only by the simulator/evaluator.
Before fitting most methods, a particle filter is run over the public dataset to
create a belief cache:
\[
  \phi_t=\phi(b_t),\qquad \phi_{t+1}=\phi(b_{t+1}),
\]
where $\phi$ contains posterior moments and quantiles of the public belief.
Current filter modes are \code{learned}, \code{raw}, \code{reference},
\code{ricker}, and \code{true\_family}; \code{oracle} is reserved for
evaluator-only ablations and is not accepted as a training input.  Cached
datasets are rejected if the population, family, reward mode, safety penalty
mode, collapse penalty, safety floor, MVP floor, or other reward-affecting cell
fields do not match the requested row.

\subsection{Train/Holdout Monitoring}

When training monitoring is enabled, dataset episodes are split into a training
set and a holdout set.  Methods fit on the training split.  The holdout split is
used to compute validation diagnostics such as dynamics error, Bellman error, or
rollout objective.  Iterative methods additionally write per-iteration curves.

\section{Shared Particle MPC}

Several methods use the same sequence-sampling planner.  Given belief $b_t$ and
a transition proposal $\hat{P}$, the planner samples candidate sequences:
\[
  A^{(j)}=(a^{(j)}_0,\ldots,a^{(j)}_{H-1}),\qquad j=1,\ldots,M.
\]
For each sequence it samples particles from $b_t$, rolls them forward under
$\hat{P}$, and computes discounted return:
\[
  G_i(A^{(j)})=\sum_{h=0}^{H-1}\gamma^h
  \hat{R}(\hat{s}^{(i)}_{t+h},a^{(j)}_h,\hat{s}^{(i)}_{t+h+1}).
\]
The score is:
\[
\begin{aligned}
  \mathrm{score}(A^{(j)})
  =&\ \frac{1}{N}\sum_iG_i(A^{(j)})\\
  &-\lambda_{\mathrm{pess}}\sqrt{\frac{1}{N}\sum_i
      \left(G_i(A^{(j)})-\bar{G}(A^{(j)})\right)^2}\\
  &-\text{optional risk penalties}\\
  &-\text{optional normalized model-disagreement penalty}.
\end{aligned}
\]
The executed action is the first action of the best sequence.

This planner is a benchmark adaptation: it gives several algorithms a common
continuous-action-evaluation engine while preserving their differences in model,
posterior, uncertainty, and constraint construction.

\section{MOPO}

\subsection{Original Paper-Level Principle}

MOPO-style offline model-based RL learns a dynamics model from offline data and
penalizes policy optimization in regions where the model is uncertain.  The core
ingredients are:
\[
  \hat{P}_\theta(s'|s,a),\qquad
  U_\theta(s,a),\qquad
  \hat{r}(s,a)-\lambda U_\theta(s,a).
\]

\subsection{Benchmark Training}

The benchmark adaptation fits a bootstrap ensemble of $B$ dynamics models.  For
each member $m$:
\[
  y_t=\log\left(1+\frac{\bar{s}_{t+1}}{K_{\mathrm{ref}}}\right),
\]
where $\bar{s}_{t+1}$ is the next-state posterior mean from the public belief
cache.  The feature vector contains log abundance features, action one-hot
features, action-state interactions, and public-control terms:
\[
\begin{aligned}
  \ell_t &= \log(1+\bar{s}_t/K_{\mathrm{ref}}),\\
  x_t &=
  \left[
  1,\ \ell_t,\ \ell_t^2,\
  e(a_t),\ e(a_t)\ell_t,\ e(a_t)\ell_t^2,\right.\\
  &\qquad\left.
  \rho_{t+1},\kappa_{t+1}/K_{\mathrm{ref}},
  K^{\mathrm{eff}}_{t+1}/K_{\mathrm{ref}},\
  e(a_t)\kappa_{t+1}/K_{\mathrm{ref}}
  \right].
\end{aligned}
\]
Each member solves a ridge regression on a bootstrap sample:
\[
  \theta_m =
  \arg\min_\theta \sum_{t\in\mathcal{B}_m}(y_t-x_t^\top\theta)^2
  +\lambda_{\mathrm{ridge}}\|\theta\|_2^2.
\]
The predicted next abundance is:
\[
  \hat{s}_{t+1}^{(m)}
  =
  K_{\mathrm{ref}}\left(\exp(x_t^\top\theta_m)-1\right).
\]

\subsection{Decision Rule}

At evaluation time, MOPO uses particle MPC with the learned ensemble as the
transition proposal.  The uncertainty signal is ensemble disagreement:
\[
  U(s,a)=\Var_m[\hat{s}^{(m)}_{t+1}(s,a)].
\]
The planner score subtracts a pessimism-scaled disagreement penalty.

\subsection{Adaptation from the Paper}

The original MOPO family is often implemented with neural dynamics and synthetic
rollout policy optimization.  The benchmark uses a ridge-regression ensemble and
particle MPC instead.  This was done because the ecological benchmark is small,
continuous-state, carefully checked for leakage, and needs deterministic,
inspectable model fits.  The adaptation preserves the core MOPO idea:
model-based planning with uncertainty pessimism, but simplifies the model class
and policy improvement machinery.

\section{Reflect-then-Plan}

\subsection{Original Paper-Level Principle}

Reflect-then-Plan methods maintain uncertainty over models or hypotheses,
update that uncertainty from observations, and plan under the posterior.  A
generic form is:
\[
  p_t(m)\propto p(o_{\le t}|m,a_{<t})p_0(m),
\]
then choose actions using posterior predictive value:
\[
  Q(b_t,a)=\E_{m\sim p_t(m)}[Q_m(b_t,a)].
\]

\subsection{Benchmark Training}

The benchmark fits the same bootstrap dynamics ensemble as MOPO, with at least
five members:
\[
  \{\hat{P}_{\theta_1},\ldots,\hat{P}_{\theta_B}\}.
\]
The initial posterior is uniform:
\[
  p_0(m)=1/B.
\]

\subsection{Decision Rule}

For a candidate sequence $A$, each selected ensemble member $m$ receives a
particle-MPC sequence score $S_m(A)$.  The reflected score is:
\[
  S(A)
  =
  \sum_m p_t(m)S_m(A)
  -
  \lambda_{\mathrm{pess}}
  \sqrt{\sum_m p_t(m)\left(S_m(A)-\sum_\ell p_t(\ell)S_\ell(A)\right)^2}.
\]
The action is the first action of the best sequence.

\subsection{Posterior Update}

After action $a_t$ and observation $o_{t+1}$, each model predicts a next-state
mean $\hat{s}_{t+1}^{(m)}$.  The update uses a log-observation residual:
\[
  \ell_m
  =
  -\frac{1}{2}\left(
    \frac{\log(o_{t+1})-\log(\hat{s}_{t+1}^{(m)})}{\sigma_m}
  \right)^2-\log\sigma_m,
\]
then:
\[
  p_{t+1}(m)=
  \frac{p_t(m)\exp(\ell_m)}
       {\sum_\ell p_t(\ell)\exp(\ell_\ell)}.
\]

\subsection{Adaptation from the Paper}

The benchmark uses the learned linear ensemble as the hypothesis set and uses a
simple residual likelihood for posterior reflection.  Planning is done by the
shared particle MPC rather than by a paper-specific planner.  This keeps the
method comparable to MOPO while preserving the essential difference: RefPlan
keeps and updates an explicit posterior over candidate learned models.

\section{BA-MCTS}

\subsection{Original Paper-Level Principle}

Bayes-adaptive MCTS plans in a belief over states and models.  It expands a tree
over actions and simulated outcomes, balancing exploration of tree actions with
estimated value:
\[
  a=\arg\max_a
  \left[
    \hat{Q}(h,a)+c\sqrt{\frac{\log N(h)}{N(h,a)}}
  \right].
\]

\subsection{Benchmark Training}

The benchmark first fits the same learned dynamics ensemble as MOPO.  The
posterior over ensemble members is initialized uniformly:
\[
  p_0(m)=1/B.
\]

\subsection{Tree Search}

At decision time, the method samples particles from the current belief and model
members from the current posterior.  Because abundance is continuous, tree nodes
use progressive aggregation:
\[
  \text{node key} =
  (d,\ \mathrm{bucket}(\log(1+s/K_{\mathrm{ref}})),\
   \1\{s\le s_{\mathrm{safe}}\},\
   \mathrm{bucket}(\rho),\mathrm{bucket}(\kappa)).
\]
For an unvisited action, the tree explores it.  Otherwise it uses UCB:
\[
  a=\arg\max_a
  \left[
    \hat{Q}(n,a)+c\sqrt{\frac{\log(N(n)+1)}{N(n,a)}}
  \right].
\]
The simulated one-step value is:
\[
  r(s,a,s')-\lambda_{\mathrm{pess}}
  \frac{\Var_m[\hat{s}'_m(s,a)]}{\max(s^2,1)}
  +\gamma V(s').
\]

\subsection{Posterior Update}

After observing a transition, the model posterior is updated with the same
log-observation residual form as Reflect-then-Plan.  The tree simulation itself
carries the public controls $(\rho,\kappa)$; the live posterior update is a
lightweight residual update on the current belief mean, action, and next
observation.

\subsection{Adaptation from the Paper}

The adaptation uses a learned ensemble as the model posterior and a continuous
state aggregation rule instead of exact discrete histories.  This is necessary
because the benchmark state and observations are continuous.  The method keeps
the paper-level structure of Bayes-adaptive tree search, but the node
representation and model posterior are benchmark-specific approximations.

\section{MOOR-Ricker}

\subsection{Original Paper-Level Principle}

MOOR-Ricker is treated as a mechanistic expert baseline: fit a simple Ricker
ecological model and plan under that model.  Its value comes from interpretably
using a known ecological structure, not from flexible model learning.

\subsection{Benchmark Training}

The method fits one Ricker model to public belief means.  The prediction is:
\[
  \hat{s}_{t+1}(r,K)
  =
  f_{\mathrm{Ricker}}(\bar{s}_t,a_t;r,K,\rho_t,\kappa_t),
\]
where the same management-control convention as the simulator is used.  The fit
objective is:
\[
  (r^\star,K^\star)
  =
  \arg\min_{r,K}
  \frac{1}{n}\sum_t
  \left[
    \log\left(1+\frac{\hat{s}_{t+1}(r,K)}{K_{\mathrm{ref}}}\right)
    -
    \log\left(1+\frac{\bar{s}_{t+1}}{K_{\mathrm{ref}}}\right)
  \right]^2.
\]
In the real set-point setting, the action-specific growth set-point is known
from the data, so the practical fit mainly identifies the capacity scale $K$.

\subsection{Decision Rule}

For real cumulative controls, MOOR plans using particle MPC with the fitted
Ricker proposal:
\[
  a_t=\pi_{\mathrm{MPC}}(b_t;\hat{P}_{\mathrm{Ricker}}(r^\star,K^\star)).
\]

\subsection{Adaptation from the Paper}

The adaptation keeps MOOR intentionally Ricker-structured even when the true
family is Allee, theta, or regime-switching.  This misspecification is part of
the baseline.  The code also uses the simulator's positive-growth/negative-
mortality split so that negative real-data set-points behave as mortality rather
than causing Ricker sign pathologies.

\section{PLUS}

\subsection{Original Paper-Level Principle}

PLUS-style methods maintain a discrete set of candidate mechanistic models,
update a posterior over those candidates from evidence, and choose actions by
planning under the posterior mixture.

\subsection{Candidate Models}

The benchmark uses $K_{\mathrm{cand}}=21$ candidates:
\[
  \mathcal{M}=\{M_1,\ldots,M_{21}\}.
\]
In the real set-point setting, the candidate axis is carrying capacity:
\[
  K_k \in [K_{\mathrm{base}},K_{\max}]
\]
with geometric spacing.  This differs from older settings where candidates may
span intrinsic growth $r$.  The reason is that real action set-points already
provide action-specific growth/mortality, while capacity uncertainty remains a
meaningful Ricker-structured uncertainty axis.

\subsection{Posterior and Belief Bank}

PLUS keeps one belief-bank copy per candidate:
\[
  b_t^{(k)},\qquad p_t(k).
\]
The initial posterior is uniform:
\[
  p_0(k)=1/K_{\mathrm{cand}}.
\]

\subsection{Decision Rule}

For a sequence $A$, each candidate produces a particle-MPC score $S_k(A)$.
The mixture score is:
\[
  S(A)=\sum_k p_t(k)S_k(A).
\]
The method executes the first action of the best sequence.

\subsection{Posterior Update}

After action $a_t$ and observation $o_{t+1}$, each candidate belief is
propagated:
\[
  b_{t+1}^{(k)}=\mathrm{BayesFilter}(b_t^{(k)},a_t,o_{t+1};M_k).
\]
The candidate evidence is:
\[
  Z_k = p(o_{t+1}|b_t^{(k)},a_t,M_k).
\]
The posterior update is:
\[
  p_{t+1}(k)=
  \frac{p_t(k)Z_k}{\sum_jp_t(j)Z_j}.
\]

\subsection{Adaptation from the Paper}

The 21-candidate count is preserved.  The adaptation changes the candidate
dimension to capacity in the real set-point setting because the real data
already specify per-action growth set-points.  PLUS remains deliberately
Ricker-assumption-based, so it serves as a structured but potentially
misspecified mechanistic baseline.

\section{Delphic-CQL}

\subsection{Original Paper-Level Principle}

Delphic-style methods consider multiple compatible latent worlds that explain
the same observations but imply different counterfactual consequences.  CQL
then learns a conservative Q-function that avoids unsupported or ambiguous
actions.

\subsection{Compatible Worlds}

The benchmark creates $W$ compatible worlds.  Given belief features $\phi_t$,
world $w$ samples a random projection $P_w$ and latent dimension $d_w$.  The
latent feature is:
\[
  z_w(\phi_t)
  =
  \tanh(\phi_tP_w)\,
  A(\phi_t)\,
  c_w,
\]
where $A(\phi_t)$ is a posterior-ambiguity scalar derived from belief spread,
and $c_w$ is a world-specific scale.  The world feature vector is:
\[
  \psi_w(\phi_t)=
  [1,\phi_t,z_w(\phi_t)].
\]

\subsection{World Fitting}

Each world fits a behavior classifier:
\[
  \pi_{\beta_w}(a|\phi)
  =
  \mathrm{softmax}(\psi_w(\phi)^\top\beta_w)_a,
\]
by gradient descent on cross-entropy.  Each world also fits a linear Q-head by
iterative Bellman regression:
\[
  y_t = R_t+\gamma(1-d_t)\max_{a'}Q_w(\phi_{t+1},a'),
\]
\[
  Q_w(\phi_t,a_t)=\psi_w(\phi_t)^\top q_{w,a_t}.
\]

\subsection{Delphic Uncertainty}

For a target action $a$, the worlds induce a set of counterfactual values.  The
Delphic uncertainty is the across-world variance:
\[
  U(\phi,a)=\Var_w[Q_w^{\mathrm{cf}}(\phi,a)].
\]

\subsection{Final CQL Head}

The final policy learns linear Q-weights $q_a$ on public belief features:
\[
  Q(\phi,a)=[1,\phi]^\top q_a.
\]
At CQL iteration $k$, the target action is
\[
  a^\star_t=\arg\max_a Q_k(\phi_t,a),
\]
and the target is:
\[
  y_t =
  R_t+\gamma(1-d_t)\max_{a'}Q_k(\phi_{t+1},a')
  -\lambda_{\mathrm{delphic}}U(\phi_t,a^\star_t).
\]
For observed actions, the head is fit by ridge regression.  Conservative
pseudo-targets are added for unsupported actions:
\[
  \min_q
  \sum_{t:a_t=a}(Q(\phi_t,a)-y_t)^2
  +\alpha_{\mathrm{CQL}}\sum_t Q(\phi_t,a)^2
  +\lambda_{\mathrm{ridge}}\|q_a\|_2^2.
\]

\subsection{Decision Rule}

At evaluation time:
\[
  a_t =
  \arg\max_a
  \left[
    Q(\phi_t,a)-\lambda_{\mathrm{delphic}}U(\phi_t,a)
  \right].
\]

\subsection{Adaptation from the Paper}

The original paper family may use richer latent world models and neural Q
functions.  The benchmark uses random-feature compatible worlds and linear CQL
heads for interpretability, speed, and leakage control.  The adaptation preserves
the core Delphic idea: penalize actions whose values vary across compatible
latent explanations of the same public evidence.

\section{OGSRL}

\subsection{Original Paper-Level Principle}

OGSRL-style methods learn a policy from offline data while guarding against
unsafe or out-of-distribution decisions.  They combine reward optimization with
safety and support constraints.

\subsection{Learned Components}

The benchmark fits:
\begin{itemize}[leftmargin=*]
  \item a learned dynamics ensemble $\hat{P}_{\theta_m}$;
  \item a KNN OOD guardian $G_{\mathrm{ood}}(s,a)$;
  \item a linear softmax actor $\pi_\eta(a|\phi)$;
  \item linear diagnostic critics for reward, safety, and OOD.
\end{itemize}

The actor is:
\[
  \pi_\eta(a|\phi)=
  \frac{\exp(\eta_a^\top \phi)}
       {\sum_b\exp(\eta_b^\top \phi)}.
\]

\subsection{OOD Guardian}

The guardian embeds offline belief/action points:
\[
  u_t=\left[
  \log(1+\bar{s}_t/K_{\mathrm{ref}}),\
  \1\{\bar{s}_t\le s_{\mathrm{safe}}\},\
  \frac{\bar{s}_t-s_{\mathrm{safe}}}{K_{\mathrm{ref}}},\
  e(a_t),\
  \rho_t,\ \kappa_t/K_{\mathrm{ref}},\
  K^{\mathrm{eff}}_t/K_{\mathrm{ref}}
  \right].
\]
It stores a subset of standardized anchors and sets a distance threshold from a
validation quantile.  A new action is OOD if its KNN distance exceeds the
threshold:
\[
  G_{\mathrm{ood}}(s,a)=
  \1\{d_k(u(s,a),\mathcal{A})>\tau\}.
\]

\subsection{Rollout Objective}

Training starts are sampled from cached belief features.  The policy is rolled
out through the learned dynamics ensemble.  At each simulated step:
\[
  r_h =
  \alpha\frac{\hat{s}_{h+1}}{\hat{s}_{h+1}+K_{\mathrm{ref}}}
  -c(a_h)
  -P_m\chi_h,
\]
and the safety/OOD costs are:
\[
  c^{\mathrm{safety}}_h=\chi_h,\qquad
  c^{\mathrm{ood}}_h=G_{\mathrm{ood}}(\hat{s}_h,a_h).
\]
Discounted returns are:
\[
  G_R=\sum_h\gamma^h r_h,\quad
  G_S=\sum_h\gamma^h c^{\mathrm{safety}}_h,\quad
  G_O=\sum_h\gamma^h c^{\mathrm{ood}}_h.
\]
The constrained actor objective is:
\[
  J(\eta)=
  \E_{\pi_\eta}\left[
    G_R-\lambda_SG_S-\lambda_OG_O
  \right].
\]

\subsection{Actor Update}

The implementation uses a REINFORCE-style score-function update.  For sampled
action $a_i$ at feature $\phi_i$:
\[
  \nabla_\eta\log\pi_\eta(a_i|\phi_i)
  =
  (\mathbf{e}_{a_i}-\pi_\eta(\cdot|\phi_i))\phi_i^\top.
\]
The update is:
\[
  \eta \leftarrow
  \eta+\alpha_\eta
  \frac{1}{N}\sum_i
  \tilde{J}_i
  \nabla_\eta\log\pi_\eta(a_i|\phi_i),
\]
where $\tilde{J}_i$ is the centered rollout objective.

Dual variables are updated by:
\[
  \lambda_S\leftarrow
  \max\{0,\lambda_S+\alpha_\lambda(\bar{G}_S-B_S)\},
\]
\[
  \lambda_O\leftarrow
  \max\{0,\lambda_O+\alpha_\lambda(\bar{G}_O-B_O)\}.
\]

\subsection{Deployment Rule}

For each action, OGSRL estimates posterior OOD probability and next-step
unsafe-occupancy probability under the current belief.  By default the training
and deployment budgets are $B_S=0.02$ and $B_O=0.05$, unless overridden by the
caller.  Actions satisfying both deployment limits are feasible.  If at least
one action is feasible:
\[
  a_t=\arg\max_{a\in\mathcal{F}_t}\pi_\eta(a|b_t).
\]
If no action is feasible, the method selects the least-violating action:
\[
  a_t=\arg\min_a
  \left[
    \frac{(\hat{p}_{\mathrm{unsafe}}(a)-B_S)_+}{B_S}
    +
    \frac{(\hat{p}_{\mathrm{ood}}(a)-B_O)_+}{B_O}
  \right].
\]

\subsection{Adaptation from the Paper}

The original paper family may use neural policies and richer OOD/safety
estimators.  The benchmark uses a linear actor, KNN support estimate with
safety/control features, and learned linear dynamics ensemble.  This makes the
method lightweight and transparent in a small ecological benchmark while
retaining the core structure: offline policy learning with safety and support
constraints.

\section{Training and Evaluation Flow}

For a single manifest row, the intended flow is:
\begin{enumerate}[leftmargin=*]
  \item instantiate one population/family/noise/reward-mode cell;
  \item resolve the requested compute backend and record the effective backend
  and fallback metadata;
  \item load or collect a public offline dataset;
  \item run the selected public filter and cache belief features;
  \item split episodes into training and holdout sets if monitoring is enabled;
  \item fit the policy on the training split;
  \item log train/holdout diagnostics;
  \item evaluate the trained policy over fixed evaluation seeds;
  \item save episode metrics, summary metrics, runtime, backend metadata, and
  training artifacts.
\end{enumerate}

The following files are intended outputs for each row:
\[
\begin{array}{ll}
\text{episodes.csv} & \text{per-episode reward, safety, abundance, filter, runtime metrics},\\
\text{summary.json} & \text{aggregate metrics and method diagnostics},\\
\text{training\_history.csv} & \text{long-form training/validation curve},\\
\text{training\_history.json} & \text{training metadata and history},\\
\text{training\_history.png} & \text{plot of loss-like curves, when enabled}.
\end{array}
\]
For real cells, the row output path is namespaced by effective backend and reward
mode, e.g. \code{backend\_numpy/reward\_safe/...}, so CPU/GPU and yield/safe
rows do not overwrite each other.

\section{Cross-Method Adaptation Summary}

\begin{center}
\begin{longtable}{p{0.16\linewidth}p{0.34\linewidth}p{0.38\linewidth}}
\toprule
Method & Paper-level core retained & Benchmark-specific adaptation \\
\midrule
\endhead
MOPO & learned dynamics plus uncertainty pessimism & ridge bootstrap ensemble plus particle MPC \\
RefPlan & posterior over models plus planning & learned ensemble posterior with residual-likelihood update \\
BA-MCTS & belief/model tree search & continuous state bucket aggregation and learned ensemble posterior \\
MOOR & simple mechanistic Ricker baseline & fit Ricker to public belief means; intentionally misspecified outside Ricker \\
PLUS & candidate mechanistic models with posterior update & 21 Ricker-style candidates over $K$ in real set-point setting \\
Delphic-CQL & compatible worlds plus conservative Q learning & random-feature latent worlds and linear CQL head \\
OGSRL & offline policy with safety/support constraints & linear actor, KNN OOD guardian with safety/control features, learned linear dynamics rollouts \\
\bottomrule
\end{longtable}
\end{center}

\section{Why These Adaptations Were Needed}

The real-ecology benchmark differs from many original algorithm papers in four
ways:
\begin{enumerate}[leftmargin=*]
  \item The state is continuous abundance, but only noisy log-normal surveys are
  observed.
  \item Actions are real management interventions with set-point mortality,
  cumulative habitat capacity, and translocation.
  \item The benchmark must preserve a strict public/private boundary.
  \item The scientific goal is controlled comparison across ecological stress
  families, not maximizing performance with large neural architectures.
\end{enumerate}

For these reasons, several methods are implemented with transparent linear models,
particle beliefs, shared particle MPC, and explicit training/holdout diagnostics.
These choices make the methods easier to inspect and compare, but they should be
interpreted as benchmark adaptations rather than exact source-code
reimplementations of every original paper.

\end{document}
